% !TeX root = closed-systems-from-comparison-completeness.tex
% ============================================================
\chapter{Diagonal Redundancy and the Obstruction to Subsystem Attribution}
\label{ch:rotation}
% ============================================================
\begin{chapterorientation}
\orientationitem{Input}{The canonical product data and diagonal group action obtained from comparison completeness.}
\orientationitem{Role}{This chapter fixes quotient semantics and proves the obstruction to assigning subsystem motion before quotient descent.}
\orientationitem{Output}{The physical quotient, orbit-equivalent twin descriptions, and the two loci through which representative-level structure may re-enter.}
\end{chapterorientation}
\section{Introduction}
Under the standing principle of closed-world admissibility
(\cref{sp:closed-world}), this chapter develops the
quotient-descent and transport-visibility clauses
(\cref{sp:quotient,sp:transport}).
Taking the canonical product data of \cref{ch:bottom} as input, it
establishes quotient descent and the subsystem-attribution boundary used
in \cref{ch:closure,ch:locus}.
\Cref{ch:bottom} established that rectangular comparison completeness
determines a canonical product presentation
\[
	X:=X_A\times X_B,
	\qquad
	X_A:=U/\alpha,
	\qquad
	X_B:=U/\beta,
\]
together with a canonical diagonal action of the intrinsic symmetry group
\[
	G:=\Aut(U,\mathcal C)
\]
on \(X\); see \cref{thm:comparison-completeness-rigidity,thm:diagonal}.
Accordingly, the primitive comparison world determines canonically the
orbit projection
\[
	\pi:X\to \Phys:=X/G,
\]
which is exactly the quotient data produced in
\cref{cor:terminal-data}.
At this stage the remaining question is semantic. Once the comparison
world has canonically collapsed to the diagonal-redundancy data
\[
	X=X_A\times X_B,
	\qquad
	G\curvearrowright X,
	\qquad
	\pi:X\to\Phys,
\]
it must be determined which state reports survive as intrinsic content
of the closed system, and where any further non-quotient structure could
possibly enter.
The first point is descent. Once the diagonal action is fixed, the
natural admissibility criterion for closed-system state reports is orbit
invariance: admissible reports are exactly the \(G\)-invariant maps on
\(X\), equivalently the maps that factor uniquely through the orbit
projection
\[
	\pi:X\to\Phys.
\]
\Cref{thm:descent} shows that this factorization is automatic.
Closed-system semantics is therefore determined to be quotient semantics.
The second point is an obstruction theorem. Within quotient semantics,
pure orbit motion carries no invariant subsystem attribution. Every
coherent report is constant along a pure diagonal orbit loop, and any
two pointwise orbit-equivalent supported representatives of such a loop
have identical coherent report values. In particular, the distinction
\begin{quote}
	which subsystem moved?
\end{quote}
does not define a quotient-level invariant of the closed system.
The third point is structural classification. Finite comparison
protocols may extend quotient semantics, but within the extension
formalism introduced below the possibilities are exhausted: if such an
extension is not endpoint-determined, then additional structure enters
through at least one of two classified loci, and possibly both:
\begin{enumerate}[label=\textup{(\roman*)}]
	\item representative selection, that is, a section of \(\pi\);
	\item morphism-level enrichment, equivalently route-dependent transport
	      or nontrivial loop defect.
\end{enumerate}
No third locus exists within that formalism.
The chapter therefore identifies the precise obstruction boundary of
closed quotient semantics. As shown below, every coherent report factors
uniquely through the orbit projection. Consequently, any invariant
content expressible by coherent reports is orbit-level content in
\(\Phys\). Within the extension formalism introduced below, any content
beyond that must arise through at least one of the two enrichment
mechanisms above, and possibly through both together.
% ============================================================
\section{Closed-system quotient semantics}
\label{sec:closed-quotient}
% ============================================================
Throughout this chapter, admissibility is read through
\cref{sp:quotient}: semantic content is required to descend through
intrinsic quotient maps, while non-quotient residue appears only through
the transport mechanisms analyzed later (\cref{sp:transport}).
Throughout the chapter, \(G\) acts diagonally on
\[
	X=X_A\times X_B,
	\qquad
	g\cdot(x_A,x_B):=(g\cdot x_A,g\cdot x_B),
\]
and
\[
	\pi:X\to \Phys:=X/G
\]
denotes the orbit projection.
\begin{definition}[Coherent report]
	\label{def:report-coherence}
	Let \(S\) be a set. A map
	\[
		R:X\to S
	\]
	is \emph{coherent} if
	\[
		R(g\cdot x)=R(x)
		\qquad
		\text{for all }g\in G,\ x\in X.
	\]
	Equivalently, \(R\) is constant on \(G\)-orbits.
\end{definition}
Because diagonal redundancy identifies orbit-equivalent representatives,
orbit invariance is not an additional convention but the intrinsic
admissibility criterion for closed-system state reports.
\begin{definition}[Closed-system semantics]
	\label{def:closed-system-semantics}
	A \emph{closed-system semantics} on \(X\) declares the admissible
	state reports to be exactly the coherent reports.
\end{definition}
\begin{definition}[Orbit quotient]
	The \emph{orbit quotient} of the diagonal action is
	\[
		\Phys:=X/G,
	\]
	equipped with the canonical projection
	\[
		\pi:X\to\Phys,
		\qquad
		\pi(x):=[x].
	\]
\end{definition}
\begin{theorem}[Canonical descent]
	\label{thm:descent}
	Let \(R:X\to S\) be coherent. Then there exists a unique map
	\[
		\widetilde R:\Phys\to S
	\]
	such that
	\[
		R=\widetilde R\circ \pi.
	\]
\end{theorem}
\begin{proof}
	Define
	\[
		\widetilde R([x]):=R(x).
	\]
	To prove that this is well defined, suppose \([x]=[y]\) in \(\Phys\).
	Then \(y=g\cdot x\) for some \(g\in G\). Since \(R\) is coherent,
	\[
		R(y)=R(g\cdot x)=R(x).
	\]
	Hence \(\widetilde R([x])\) is independent of the chosen representative.
	Now for every \(x\in X\),
	\[
		(\widetilde R\circ\pi)(x)
		=
		\widetilde R([x])
		=
		R(x),
	\]
	so \(R=\widetilde R\circ\pi\).
	For uniqueness, let \(f:\Phys\to S\) also satisfy \(R=f\circ\pi\). Then
	for any \([x]\in\Phys\),
	\[
		f([x])=f(\pi(x))=R(x)=\widetilde R([x]).
	\]
	Therefore \(f=\widetilde R\).
\end{proof}
\begin{proposition}[Universal property of the orbit quotient]
	\label{prop:orbit-universal}
	For every set \(S\) and every coherent map
	\[
		R:X\to S,
	\]
	there exists a unique map
	\[
		\widetilde R:\Phys\to S
	\]
	such that
	\[
		R=\widetilde R\circ\pi.
	\]
	Equivalently, \(\pi:X\to\Phys\) is universal among coherent maps out of \(X\).
\end{proposition}
\begin{proof}
	This is exactly \cref{thm:descent}.
\end{proof}
\begin{corollary}[Coherence \(\Longleftrightarrow\) descent]
	\label{cor:descent}
	For a map \(R:X\to S\), the following are equivalent.
	\begin{enumerate}[label=\textup{(\roman*)}]
		\item \(R\) is coherent;
		\item \(R\) is constant on \(G\)-orbits;
		\item there exists a unique \(\widetilde R:\Phys\to S\) such that
		      \[
			      R=\widetilde R\circ\pi.
		      \]
	\end{enumerate}
\end{corollary}
\begin{proof}
	\emph{(i)\(\Rightarrow\)(iii).}
	This is \cref{thm:descent}.
	\smallskip
	\emph{(iii)\(\Rightarrow\)(ii).}
	If \(R=\widetilde R\circ\pi\), then for \(x,y\in X\) with
	\(\pi(x)=\pi(y)\),
	\[
		R(x)=\widetilde R(\pi(x))=\widetilde R(\pi(y))=R(y).
	\]
	Thus \(R\) is constant on orbits.
	\smallskip
	\emph{(ii)\(\Rightarrow\)(i).}
	If \(R\) is constant on orbits, then for every \(g\in G\) and
	\(x\in X\), the points \(x\) and \(g\cdot x\) lie in the same orbit, so
	\[
		R(g\cdot x)=R(x).
	\]
	Hence \(R\) is coherent.
\end{proof}
\begin{remark}[Structural meaning of quotient semantics]
	The orbit projection
	\[
		\pi:X\to \Phys=X/G
	\]
	is not merely a convenient quotient map. By
	\cref{prop:orbit-universal}, it is universal among coherent
	descriptions of \(X\). Thus quotient semantics is determined by diagonal
	redundancy: every admissible state report factors uniquely through
	\(\Phys\), and no finer representative-level distinction survives
	without additional structure.
\end{remark}
\begin{remark}[Categorical formulation]
	Let \(\mathsf{Inv}(X)\) be the category whose objects are pairs
	\((S,R)\), where \(S\) is a set and \(R:X\to S\) is coherent, and
	whose morphisms
	\[
		f:(S,R)\to (S',R')
	\]
	are maps \(f:S\to S'\) satisfying
	\[
		R'=f\circ R.
	\]
	Then \((\Phys,\pi)\) is an initial object of \(\mathsf{Inv}(X)\).
	This is precisely the categorical form of
	\cref{prop:orbit-universal}.
\end{remark}
% ============================================================
\section{Pure orbit loops}
\label{sec:loops}
% ============================================================
We now analyze histories that move entirely inside a single diagonal
orbit.
\begin{definition}[Pure diagonal orbit loop]
	Fix \(T>0\). A \emph{pure diagonal orbit loop} is a map
	\[
		\Gamma:[0,T]\to X
	\]
	for which there exists a map
	\[
		g:[0,T]\to G
	\]
	such that
	\[
		g(0)=g(T)=e
		\qquad\text{and}\qquad
		\Gamma(t)=g(t)\cdot \Gamma(0)
		\quad
		\text{for all }t\in[0,T].
	\]
\end{definition}
\begin{definition}[\(A\)-supported and \(B\)-supported representatives]
	Let \(\Gamma:[0,T]\to X\) be a path.
	An \emph{\(A\)-supported representative} of \(\Gamma\) is a path
	\[
		\Gamma_A(t)=(x_A(t),x_B(0))
	\]
	such that
	\[
		\pi(\Gamma_A(t))=\pi(\Gamma(t))
		\qquad
		\text{for all }t\in[0,T].
	\]
	Similarly, a \emph{\(B\)-supported representative} is a path
	\[
		\Gamma_B(t)=(x_A(0),x_B(t))
	\]
	satisfying
	\[
		\pi(\Gamma_B(t))=\pi(\Gamma(t))
		\qquad
		\text{for all }t\in[0,T].
	\]
\end{definition}
\begin{lemma}[Single-orbit property]
	\label{lem:single-orbit}
	If \(\Gamma\) is a pure diagonal orbit loop, then
	\[
		\pi(\Gamma(t))=\pi(\Gamma(0))
		\qquad
		\text{for all }t\in[0,T].
	\]
\end{lemma}
\begin{proof}
	For each \(t\in[0,T]\), the defining relation
	\[
		\Gamma(t)=g(t)\cdot \Gamma(0)
	\]
	shows that \(\Gamma(t)\) lies in the same \(G\)-orbit as \(\Gamma(0)\).
	Since \(\pi\) is the orbit projection, it takes equal values on
	orbit-equivalent points. Therefore
	\[
		\pi(\Gamma(t))=\pi(\Gamma(0))
	\]
	for all \(t\).
\end{proof}
\begin{remark}
	A pure orbit loop may be descriptively nontrivial in the product
	presentation \(X_A\times X_B\), but orbit-theoretically it is
	trivial: it remains inside a single orbit of the diagonal action.
\end{remark}
% ============================================================
\section{Orbit-equivalent supported representatives}
\label{sec:twin}
% ============================================================
The following proposition isolates the elementary symmetry relating
supported representatives of the same quotient path.
\begin{proposition}[Orbit-equivalent supported representatives]
	\label{prop:twin}
	Let \(\Gamma:[0,T]\to X\) be a pure diagonal orbit loop with witness
	\(g:[0,T]\to G\), and assume \(\Gamma\) admits an \(A\)-supported
	representative \(\Gamma_A\). Define
	\[
		\widetilde\Gamma(t):=g(t)^{-1}\cdot \Gamma_A(t).
	\]
	Then:
	\begin{enumerate}[label=\textup{(\roman*)}]
		\item \(\widetilde\Gamma\) is a path in \(X\) satisfying
		      \[
			      \pi(\widetilde\Gamma(t))=\pi(\Gamma(t))
			      \qquad
			      \text{for all }t\in[0,T];
		      \]
		\item one has the reconstruction identity
		      \[
			      \Gamma_A(t)=g(t)\cdot \widetilde\Gamma(t)
			      \qquad
			      \text{for all }t\in[0,T];
		      \]
		\item the property of admitting an \(A\)-supported representative
		      depends only on the quotient path \(t\mapsto \pi(\Gamma(t))\);
		\item if a \(B\)-supported representative of the same quotient path is
		      chosen, then both \(\Gamma\) and \(\widetilde\Gamma\) are pointwise
		      orbit-equivalent to it.
	\end{enumerate}
\end{proposition}
\begin{proof}
	For each \(t\in[0,T]\),
	\[
		\widetilde\Gamma(t)=g(t)^{-1}\cdot\Gamma_A(t)
	\]
	lies in the same \(G\)-orbit as \(\Gamma_A(t)\). Since \(\Gamma_A\) is a
	representative of \(\Gamma\), we have
	\[
		\pi(\Gamma_A(t))=\pi(\Gamma(t)).
	\]
	Therefore
	\[
		\pi(\widetilde\Gamma(t))=\pi(\Gamma_A(t))=\pi(\Gamma(t)),
	\]
	which proves \textup{(i)}.
	Statement \textup{(ii)} is immediate from the definition:
	\[
		g(t)\cdot \widetilde\Gamma(t)
		=
		g(t)\cdot \bigl(g(t)^{-1}\cdot \Gamma_A(t)\bigr)
		=
		\Gamma_A(t).
	\]
	For \textup{(iii)}, the existence of an \(A\)-supported representative
	is, by definition, a statement about whether the quotient path
	\(t\mapsto \pi(\Gamma(t))\) can be realized by a path of the form
	\((x_A(t),x_B(0))\). Thus it depends only on the quotient path itself.
	For \textup{(iv)}, let \(\Gamma_B\) be a \(B\)-supported representative
	of the same quotient path. Then
	\[
		\pi(\Gamma_B(t))=\pi(\Gamma(t))=\pi(\widetilde\Gamma(t))
		\qquad
		\text{for all }t,
	\]
	so \(\Gamma_B(t)\), \(\Gamma(t)\), and \(\widetilde\Gamma(t)\) are
	pairwise orbit-equivalent at each time.
\end{proof}
% ============================================================
\section{Subsystem attribution obstruction}
\label{sec:nogo}
% ============================================================
We now prove the main no-go statement.
\begin{theorem}[No-go theorem]
	\label{thm:nogo}
	Let \(R:X\to S\) be coherent, and let \(\Gamma:[0,T]\to X\) be a pure
	diagonal orbit loop. Then
	\[
		R\circ \Gamma
	\]
	is constant on \([0,T]\).
	Moreover, if \(\Gamma\) admits an \(A\)-supported representative and
	\(\widetilde\Gamma\) is the associated path from \cref{prop:twin}, then
	\[
		R(\widetilde\Gamma(t))=R(\Gamma(t))
		\qquad
		\text{for all }t\in[0,T].
	\]
\end{theorem}
\begin{proof}
	By \cref{thm:descent}, there exists a unique map
	\[
		\widetilde R:\Phys\to S
	\]
	such that
	\[
		R=\widetilde R\circ\pi.
	\]
	Since \(\Gamma\) is a pure diagonal orbit loop,
	\cref{lem:single-orbit} gives
	\[
		\pi(\Gamma(t))=\pi(\Gamma(0))
		\qquad
		\text{for all }t\in[0,T].
	\]
	Hence for all \(t\),
	\[
		R(\Gamma(t))
		=
		\widetilde R(\pi(\Gamma(t)))
		=
		\widetilde R(\pi(\Gamma(0))),
	\]
	which is independent of \(t\). Thus \(R\circ\Gamma\) is constant.
	For the second statement, \cref{prop:twin}(i) gives
	\[
		\pi(\widetilde\Gamma(t))=\pi(\Gamma(t))
		\qquad
		\text{for all }t\in[0,T].
	\]
	Therefore
	\[
		R(\widetilde\Gamma(t))
		=
		\widetilde R(\pi(\widetilde\Gamma(t)))
		=
		\widetilde R(\pi(\Gamma(t)))
		=
		R(\Gamma(t)).
	\]
\end{proof}
\begin{corollary}[No subsystem attribution at quotient level]
	Within quotient semantics, subsystem attribution for pure orbit motion
	does not define an invariant: every coherent report is constant along
	pure orbit loops, and any two pointwise orbit-equivalent
	representatives have identical coherent report values.
\end{corollary}
\begin{proof}
	The constancy statement is the first part of \cref{thm:nogo}. The
	pointwise orbit-equivalence statement follows from the second part and
	\cref{prop:twin}(iv).
\end{proof}
% ============================================================
\section{Orbit-preserving maps}
\label{sec:orbit-invisibility}
% ============================================================
\Cref{thm:nogo} is the pathwise case of a more general
orbit-invisibility principle.
\begin{definition}[Orbit-preserving map]
	A map
	\[
		\Theta_X:X\to X
	\]
	is \emph{orbit-preserving} if
	\[
		\pi\circ\Theta_X=\pi.
	\]
	Equivalently, \(\Theta_X(x)\) lies in the same \(G\)-orbit as \(x\) for
	every \(x\in X\).
\end{definition}
\begin{theorem}[Orbit-preserving maps are coherently trivial]
	Let
	\[
		\Theta_X:X\to X
	\]
	be orbit-preserving. Then for every coherent report \(R:X\to S\),
	\[
		R\circ \Theta_X=R.
	\]
\end{theorem}
\begin{proof}
	By \cref{thm:descent}, write
	\[
		R=\widetilde R\circ\pi
	\]
	for a unique map \(\widetilde R:\Phys\to S\). Then for every \(x\in X\),
	\[
		R(\Theta_X(x))
		=
		\widetilde R(\pi(\Theta_X(x)))
		=
		\widetilde R(\pi(x))
		=
		R(x),
	\]
	because \(\Theta_X\) is orbit-preserving.
\end{proof}
\begin{remark}
	Within quotient semantics, an orbit-preserving map introduces no new
	invariant information. Any nontrivial distinction must therefore arise
	from structure not already encoded by the quotient projection.
\end{remark}
% ============================================================
\section{Finite protocol extensions and the two-locus classification}
\label{sec:two-locus-classification}
% ============================================================
We now classify all ways of enriching quotient semantics on finite
comparison protocols.
\begin{definition}[Finite protocol network]
	A \emph{finite protocol network} is a finite directed graph
	\[
		\mathcal N=(V,E)
	\]
	whose vertex set satisfies
	\[
		V\subseteq \Phys.
	\]
	Write \(\Path(\mathcal N)\) for the free category on \(\mathcal N\).
\end{definition}
\begin{definition}[Extension of quotient semantics]
	\label{def:extension}
	Let \(\mathcal N=(V,E)\) be a finite protocol network. An
	\emph{extension of quotient semantics} consists of the following data.
	\begin{enumerate}[label=\textup{(\roman*)}]
		\item \emph{State compatibility:} for each \(p\in V\), an object
		      assignment
		      \[
			      p\longmapsto F_{\mathcal N}(p)
		      \]
		      depending only on \(p\in\Phys\);
		\item \emph{Compositionality:} a small category
		      \[
			      \mathsf C_{\mathcal N}
		      \]
		      with object set \(V\), together with a functor
		      \[
			      F_{\mathcal N}:\Path(\mathcal N)\to\mathsf C_{\mathcal N}
		      \]
		      extending the object assignment.
	\end{enumerate}
\end{definition}
\begin{definition}[Endpoint-determined extension]
	An extension of quotient semantics is \emph{endpoint-determined} if for
	every finite protocol network \(\mathcal N\) and every pair of
	co-terminal paths
	\[
		\gamma_1,\gamma_2\in\Path(\mathcal N)
	\]
	one has
	\[
		F_{\mathcal N}(\gamma_1)=F_{\mathcal N}(\gamma_2).
	\]
\end{definition}
\begin{definition}[Section]
	A \emph{section} of the orbit projection is a map
	\[
		s:\Phys\to X
	\]
	such that
	\[
		\pi\circ s=\id_{\Phys}.
	\]
\end{definition}
\begin{theorem}[Two-locus classification within quotient-semantic extensions]
	\label{thm:two-locus-classification}
	Under the standing principle \cref{sp:closed-world}, let \(F\) be an
	\emph{extension of quotient semantics} in the sense of
	\cref{def:extension}. If \(F\) is not endpoint-determined, then at
	least one of the following two enrichment mechanisms occurs:
	\begin{enumerate}[label=\textup{(\roman*)}]
		\item \emph{Representative selection:} a section
		      \[
			      s:\Phys\to X
		      \]
		      of \(\pi\) is supplied;
		\item \emph{Morphism enrichment:} there exist a network
		      \[
			      \mathcal N=(V,E)
		      \]
		      and co-terminal paths
		      \[
			      \gamma_1,\gamma_2\in\Path(\mathcal N)
		      \]
		      such that
		      \[
			      F_{\mathcal N}(\gamma_1)\neq F_{\mathcal N}(\gamma_2).
		      \]
	\end{enumerate}
	No third locus exists within the data specified in \cref{def:extension}.
\end{theorem}
\begin{proof}
	State compatibility fixes the object-level assignment by the points
	\(p\in\Phys\) alone. Thus non-endpoint behavior cannot arise from a new
	interpretation of vertices as such.
	Accordingly, any failure of endpoint determinacy must be witnessed
	through at least one of two enrichment layers, and both may be
	present together. First, one may adjoin representative-level data
	lifting points of \(\Phys\) back to \(X\). This is exactly the supply of a section
	\[
		s:\Phys\to X,
	\]
	which chooses a distinguished representative in each orbit.
	Second, one may assign genuinely distinct morphism data to
	co-terminal paths. Since \(\Path(\mathcal N)\) is the free category
	on the graph \(\mathcal N\), a functor out of it is determined by its
	values on objects, generating edges, and composition. The object
	layer is already fixed by state compatibility. Therefore any
	remaining failure of endpoint determinacy occurs at the morphism
	layer, whether or not a section is also present.
	This proves that mechanisms \textup{(i)} and \textup{(ii)} exhaust all
	possible sources of non-endpoint behavior. No third locus exists
	because \cref{def:extension} contains no further data.
\end{proof}
\begin{remark}[Object--morphism complementarity]
	The two enrichment mechanisms act on categorically orthogonal layers.
	Sections modify the representative choice at the object level.
	Transport or holonomy data modify the morphism layer while leaving the
	object set unchanged. Under the axioms of \cref{def:extension},
	nothing else is available.
\end{remark}
% ============================================================
\section{Transport schemes and holonomy}
\label{sec:transport}
% ============================================================
We now recast the morphism-locus part of the classification as
route-dependent transport.
\begin{definition}[Free path groupoid]
	\label{def:pathpm}
	Let \(\mathcal N=(V,E)\) be a finite directed graph. Write
	\[
		\Path^\pm(\mathcal N)
	\]
	for the free groupoid obtained from \(\mathcal N\) by adjoining formal
	inverses to all directed edges, in the usual universal-property sense
	\cite{MacLane1998}.
\end{definition}
\begin{definition}[Comparison groupoid]
	A \emph{comparison groupoid} on \(\mathcal N\) is a small groupoid
	\[
		\mathsf C_{\mathcal N}\rightrightarrows V.
	\]
\end{definition}
\begin{definition}[Transport scheme]
	A \emph{transport scheme} on a finite network \(\mathcal N=(V,E)\)
	consists of a choice, for each edge \(e:p\to q\), of a morphism
	\[
		\tau_e\in \Hom_{\mathsf C_{\mathcal N}}(p,q).
	\]
	By the universal property of the free groupoid, these data extend
	uniquely to a functor
	\[
		\Hol:\Path^\pm(\mathcal N)\to \mathsf C_{\mathcal N}
	\]
	satisfying
	\[
		\Hol(e)=\tau_e
		\qquad
		\text{for every generating edge }e.
	\]
\end{definition}
\begin{definition}[Route dependence]
	A transport scheme is \emph{endpoint-determined} if for any two
	morphisms
	\[
		\gamma_1,\gamma_2:p\to q
	\]
	in \(\Path^\pm(\mathcal N)\) one has
	\[
		\Hol(\gamma_1)=\Hol(\gamma_2).
	\]
	Otherwise it is \emph{route-dependent}.
\end{definition}
\begin{lemma}[Loop reduction]
	\label{lem:loop-based}
	Let
	\[
		\gamma_1,\gamma_2:p\to q
	\]
	be morphisms in \(\Path^\pm(\mathcal N)\). Then
	\[
		\Hol(\gamma_1)=\Hol(\gamma_2)
		\iff
		\Hol(\gamma_2^{-1}\!\circ\gamma_1)=\id_p.
	\]
\end{lemma}
\begin{proof}
	Because \(\mathsf C_{\mathcal N}\) is a groupoid, \(\Hol(\gamma_2)\) is
	invertible. Therefore
	\[
		\Hol(\gamma_2^{-1}\!\circ\gamma_1)
		=
		\Hol(\gamma_2)^{-1}\circ \Hol(\gamma_1).
	\]
	Thus \(\Hol(\gamma_2^{-1}\!\circ\gamma_1)=\id_p\) if and only if
	\[
		\Hol(\gamma_1)=\Hol(\gamma_2).
	\]
\end{proof}
\begin{theorem}[Route dependence \(\Longleftrightarrow\) loop defect]
	\label{thm:route-dependence}
	A transport scheme is route-dependent if and only if there exists a
	based loop
	\[
		\ell:p\to p
	\]
	in \(\Path^\pm(\mathcal N)\) such that
	\[
		\Hol(\ell)\neq \id_p.
	\]
\end{theorem}
\begin{proof}
	Suppose first that the transport scheme is route-dependent. Then there
	exist co-terminal morphisms
	\[
		\gamma_1,\gamma_2:p\to q
	\]
	such that
	\[
		\Hol(\gamma_1)\neq \Hol(\gamma_2).
	\]
	Set
	\[
		\ell:=\gamma_2^{-1}\!\circ\gamma_1.
	\]
	Then by \cref{lem:loop-based},
	\[
		\Hol(\ell)
		=
		\Hol(\gamma_2)^{-1}\circ\Hol(\gamma_1)
		\neq \id_p.
	\]
	Conversely, suppose every based loop has trivial holonomy. Let
	\[
		\gamma_1,\gamma_2:p\to q
	\]
	be co-terminal. Then
	\[
		\gamma_2^{-1}\!\circ\gamma_1
	\]
	is a loop at \(p\), so
	\[
		\Hol(\gamma_2^{-1}\!\circ\gamma_1)=\id_p.
	\]
	By \cref{lem:loop-based},
	\[
		\Hol(\gamma_1)=\Hol(\gamma_2).
	\]
	Thus the scheme is endpoint-determined.
\end{proof}
\begin{corollary}[Transport form of the two-locus classification]
	\label{cor:transport-form}
	For finite protocol extensions in the sense of \cref{def:extension},
	any enrichment beyond quotient semantics must occur through at least
	one of the following classified loci, and may occur through both:
	\begin{enumerate}[label=\textup{(\roman*)}]
		\item object-level enrichment by a section of \(\pi\);
		\item morphism-level enrichment by a route-dependent transport scheme,
		      equivalently by nontrivial loop defect.
	\end{enumerate}
\end{corollary}
\begin{proof}
	Combine \cref{thm:two-locus-classification} with
	\cref{thm:route-dependence}.
\end{proof}
\begin{remark}[Positive interpretation]
	The two-locus classification says exactly where additional structure
	can enter once quotient semantics has been fixed. A section chooses
	preferred representatives in each orbit. A route-dependent transport
	scheme introduces morphism-level data not determined by the endpoints
	in \(\Phys\). Absent both of these two enrichments, no further
	invariant distinction exists beyond orbit-level content.
\end{remark}
% ============================================================
\section{Conclusion}
\label{sec:rotation-conclusion}
Starting from the canonical diagonal redundancy data produced in
\cref{ch:bottom},
\[
	X=X_A\times X_B,
	\qquad
	G=\Aut(U,\mathcal C),
	\qquad
	\pi:X\to\Phys:=X/G,
\]
the semantic consequences of diagonal symmetry are now fixed.
First, admissible closed-system reports are exactly the coherent maps
\[
	R:X\to S
\]
that are invariant under the diagonal action of \(G\). By the descent
theorem, every such report factors uniquely through the orbit
projection
\[
	\pi:X\to\Phys.
\]
Closed-system semantics is therefore determined to be quotient semantics:
all invariant reportable content is carried by the orbit space
\(\Phys\).
Second, pure orbit motion is invisible to coherent reports. Every
coherent report is constant along a pure diagonal orbit loop, and any
pointwise orbit-equivalent supported representatives of such a loop
have identical coherent report values. Hence subsystem attribution for
pure orbit motion does not define a quotient-level invariant.
Third, the possible enrichments of quotient semantics on finite
comparison protocols are completely classified. Within the extension
formalism of \cref{def:extension}, the classification is exhaustive: if
an extension is not endpoint-determined, then additional structure
enters through one or both of the following mechanisms:
representative selection through a section of \(\pi\), and
morphism-level transport data, equivalently nontrivial loop defect.
There is no third locus within that formalism.
The chapter therefore identifies the precise boundary of closed
quotient semantics. Diagonal redundancy collapses all intrinsic
invariant content of the closed comparison world to the orbit quotient
\[
	\Phys=X/G.
\]
Any further distinction must therefore be imported explicitly through
at least one of the two structural mechanisms above, and possibly
through both together.
Accordingly, the semantic frontier is exact: intrinsic closed-system
content is exhausted by quotient data, and every admissible refinement is
confined to the classified object and morphism loci.

\Cref{ch:closure} compresses this semantic arc into a single closure
theorem that functions as the entry statement for the subsequent
classification and obstruction chapters.

\begin{chapterclosing}
\closingitem{Established}{All coherent closed-system reports descend through \(\Phys=X/G\), and any admissible representative enrichment lies in the object or morphism locus.}
\closingitem{Carried forward}{The quotient-semantic result is consolidated next as the formal closure theorem.}
\end{chapterclosing}
